% =====================================================================
%  CARNET D'ENTRAÎNEMENT — Fiche 12 (Chimie organique)
%  THÈME 1 — Hydrocarbures, formules brutes et squelette carboné
%  NIVEAU FACILE — ÉNONCÉS
% =====================================================================
\documentclass[11pt,a4paper]{article}
\usepackage[utf8]{inputenc}
\usepackage[T1]{fontenc}
\usepackage{lmodern}
\usepackage[french]{babel}
\usepackage{amsmath,amssymb,mathtools}
\providecommand{\degree}{^\circ}
\usepackage{icomma}
\usepackage{siunitx}
\sisetup{output-decimal-marker={,},per-mode=symbol}
\usepackage[version=4]{mhchem}
\usepackage{chemfig}
\setchemfig{atom sep=2em,bond length=0.9em,cram width=2.5pt}
\usepackage{xcolor}
\usepackage{array}
\usepackage{booktabs}
\usepackage{tabularx}
\usepackage[most]{tcolorbox}
\usepackage{enumitem}
\usepackage{tikz}
\usetikzlibrary{arrows.meta,positioning,calc}
\usepackage{fancyhdr}
\usepackage{titlesec}
\usepackage[hidelinks]{hyperref}
\usepackage[a4paper,margin=2.1cm,headheight=26pt,headsep=9pt]{geometry}

\definecolor{primary}{RGB}{150,98,162}
\definecolor{primarydark}{RGB}{92,52,110}
\definecolor{excol}{RGB}{0,135,90}
\definecolor{attncol}{RGB}{200,40,55}
\definecolor{methcol}{RGB}{90,90,100}

\newcommand{\runtitle}{Thème 1 — Facile — Énoncés}
\pagestyle{fancy}\fancyhf{}
\fancyhead[L]{\small\textcolor{primarydark}{\textbf{Fiche 12} — Chimie organique}}
\fancyhead[R]{\small\textcolor{primarydark}{\runtitle}}
\fancyfoot[C]{\small\thepage}
\renewcommand{\headrulewidth}{0.8pt}
\renewcommand{\headrule}{\hbox to\headwidth{\color{primary}\leaders\hrule height \headrulewidth\hfill}}
\setlength{\parindent}{0pt}\setlength{\parskip}{3pt}

\tcbset{boxbase/.style={enhanced,breakable,boxrule=0.9pt,arc=2.5pt,
   left=3mm,right=3mm,top=2mm,bottom=2mm,fonttitle=\bfseries,coltitle=white,
   toptitle=0.5mm,bottomtitle=0.5mm}}
\newtcolorbox[auto counter]{exo}[1]{boxbase,colback=primary!4,colframe=primary,
   title={\textsc{Exercice}~\thetcbcounter\ \textmd{--- \itshape #1}}}
\newtcolorbox{consigne}{boxbase,colback=methcol!8,colframe=methcol,sharp corners,
   title={\textsc{Consignes}}}

\newcommand{\banner}[2]{%
\begin{tcolorbox}[enhanced,colback=primary,colframe=primary,arc=5pt,boxrule=0pt,
   left=5mm,right=5mm,top=2.5mm,bottom=2.5mm,coltext=white]
{\large\bfseries #1}\\[1pt]{\small #2}
\end{tcolorbox}\vspace{2pt}}

\begin{document}

\banner{Thème 1 — Hydrocarbures, formules brutes et squelette}{Niveau \textbf{Facile} — Énoncés \quad$\vert$\quad Fiche 12, Chimie organique}

\begin{consigne}
Exercices \textbf{ouverts} : rédige tes réponses (pas de QCM). Chaque exercice isole \emph{une} mécanique du tuto. Aucune calculatrice n'est nécessaire. Formules générales : alcane $\ce{C_nH_{2n+2}}$, alcène $\ce{C_nH_{2n}}$, alcyne $\ce{C_nH_{2n-2}}$.
\end{consigne}

% ---------------------------------------------------------------------
\begin{exo}{appliquer les formules générales}
Pour une chaîne acyclique de \textbf{6 atomes de carbone}, écris la formule brute :
\begin{enumerate}[label=\textbf{\alph*)},leftmargin=2em,itemsep=1pt]
  \item de l'alcane correspondant ;
  \item de l'alcène (une seule double liaison) correspondant ;
  \item de l'alcyne (une seule triple liaison) correspondant.
\end{enumerate}
Indique à chaque fois le nom de la molécule la plus simple portant ce squelette (hexane, hex-1-ène, hex-1-yne).
\end{exo}

% ---------------------------------------------------------------------
\begin{exo}{compter les hydrogènes de chaque carbone}
On donne le \textbf{2-méthylbutane}, de formule semi-développée $\ce{CH3-CH(CH3)-CH2-CH3}$ :
\begin{center}\chemfig{CH_3-[:30]CH(-[:-90]CH_3)-[:-30]CH_2-[:30]CH_3}\end{center}
En raisonnant « chaque carbone complète ses liaisons jusqu'à 4 », indique le \textbf{nombre d'atomes d'hydrogène} porté par chacun des \textbf{5} carbones. Vérifie ensuite que le total redonne la formule brute $\ce{C5H12}$.
\end{exo}

% ---------------------------------------------------------------------
\begin{exo}{classer les carbones (I / II / III / IV)}
Toujours sur le \textbf{2-méthylbutane} ci-dessus, classe chacun de ses \textbf{5 carbones} en \emph{primaire}, \emph{secondaire}, \emph{tertiaire} ou \emph{quaternaire}, en te fondant uniquement sur le nombre d'\emph{autres carbones} auxquels il est lié.
\end{exo}

% ---------------------------------------------------------------------
\begin{exo}{calculer un degré d'insaturation}
Calcule l'indice de déficience en hydrogène $I=\dfrac{2n+2-m}{2}$ pour chacune des formules brutes suivantes, puis conclus sur la nature de l'hydrocarbure (saturé acyclique, ou présence de liaisons multiples / cycle) :
\begin{center}
\ce{C4H10} \qquad \ce{C4H8} \qquad \ce{C4H6}
\end{center}
\end{exo}

% ---------------------------------------------------------------------
\begin{exo}{reconnaître la famille sur un squelette}
Voici trois molécules à \textbf{5 carbones}, en écriture topologique. Pour chacune, donne la \textbf{famille} (alcane / alcène / alcyne) et la \textbf{formule brute}.
\begin{center}
\begin{tabular}{ccc}
\chemfig{-[:30]-[:-30]-[:30]-[:-30]} &
\chemfig{-[:30]=[:-30]-[:30]-[:-30]} &
\chemfig{-[:30]~[:-30]-[:30]-[:-30]} \\[4pt]
\textbf{(A)} & \textbf{(B)} & \textbf{(C)}
\end{tabular}
\end{center}
\end{exo}

\end{document}
